\section{The parameter field detected by logarithmic forms}
\label{sec:inputs}\label{it:section}

We first find a subfield of the original base field whose transcendence
degree is exactly the excess of the total logarithmic Kodaira dimension
over the fiber Kodaira dimension. This construction uses logarithmic
subadditivity and section maps; constancy of the fibers will be proved
after that field has been identified.

\subsection{The lower bound and section conventions}

The precise lower-bound input is the following companion theorem.
Only its projective form is used here. A compact manifold belongs to
Fujiki class $\mathcal C$ if it is bimeromorphic to a compact K\"ahler
manifold; in particular, every projective manifold belongs to this class.

\begin{theorem}[Reduced-SNC logarithmic subadditivity]
\label{spec:logarithmic}
Let $f:X\to Y$ be a surjective morphism with connected fibers between
smooth connected projective complex varieties. Let $D_X,D_Y$ be
reduced simple normal crossing divisors, with the zero divisor allowed,
such that
\[
 \operatorname{Supp}(f^*D_Y)\subseteq\operatorname{Supp}(D_X).
\]
For a very general fiber $F$, put $D_F=D_X|_F$. Then
\[
 \kappa(X,K_X+D_X)
 \geq\kappa(F,K_F+D_F)+\kappa(Y,K_Y+D_Y).
\]
The same assertion holds for compact manifolds in Fujiki class
$\mathcal C$ and holomorphic fibrations between them.
\end{theorem}

\begin{proof}
Apply
\cite[\PoneLog]{SubadditivityCompanion}.
The zero-divisor and $-\infty$ conventions here agree with that statement.
Only its projective form is used in the present paper.
\end{proof}

\subsection{Models, sections, and fields}

In the algebraic arguments below, varieties are integral unless a
finite splitting algebra is explicitly considered. For a surjective morphism with connected fibers between normal
varieties in characteristic zero, Stein factorization makes the base
field relatively algebraically closed in the total function field; the
generic field extension is regular and the generic fiber is geometrically
integral.
We use smooth projective resolutions of varieties and pairs
\cite{BierstoneMilmanFunctorial}; resolving the closure of a graph
also eliminates indeterminacy. Boundaries on these models will always
be specified. A rational section of a line bundle on a normal variety
is regular exactly when it is regular at every prime divisor.

\begin{lemma}[Exceptional section comparison]
\label{prelim:sections}\label{setup:log-modification}
Let \((X,\Delta)\) be a smooth pair with effective SNC rational boundary of coefficients at most one. Let \(\mu:X'\to X\) be a smooth log resolution, and let \(\Delta'\) be the strict transform of \(\Delta\) plus the reduced exceptional divisor. For every sufficiently divisible \(m\) and every line bundle \(A\) on \(X\), pullback and pushdown identify
\[
H^0\bigl(X,m(K_X+\Delta)+A\bigr)
\simeq
H^0\bigl(X',m(K_{X'}+\Delta')+\mu^*A\bigr).
\]
More precisely, there is an effective exceptional rational divisor $E$
with $K_{X'}+\Delta'=\mu^*(K_X+\Delta)+E$, and
\[
 \mu_*\cO_{X'}\bigl(m(K_{X'}+\Delta')\bigr)
   =\cO_X\bigl(m(K_X+\Delta)\bigr).
\]
This sheaf identity respects multiplication and remains valid after
pushing to any common base. More generally, adding an effective
exceptional divisor to a pulled-back rational divisor does not change
its divisible section spaces.
\end{lemma}

\begin{proof}
The discrepancy is exceptional because a birational morphism between
smooth varieties is an isomorphism at every generic target divisor.
For an exceptional prime $E_0$, choose boundary coordinates $z_i$ on
$X$, let $a_i=\ord_{E_0}(\mu^*z_i)$, and let $k_{E_0}$ be its
Jacobian order. Pulling back a logarithmic top wedge gives
$k_{E_0}+1\geq\sum_i a_i$: a nonzero wedge contains at most one
factor $du/u$ at $E_0=(u=0)$. If the boundary coefficients are
$\delta_i\leq1$, the discrepancy coefficient is therefore
\[
 k_{E_0}+1-\sum_i\delta_i a_i
       \geq\sum_i(1-\delta_i)a_i\geq0.
\]
Thus $E$ is effective. A rational function with poles only on an
effective exceptional divisor passes every codimension-one regularity
test on the normal target, so $\mu_*\cO_{X'}(mE)=\cO_X$.
The projection formula proves the sheaf and section identities,
including their twisted, relative, and multiplicative forms.
\end{proof}

The same comparison applies after restricting to a sufficiently general smooth fiber: the finitely many relevant horizontal strata meet it transversely, and vertical exceptional loci can be avoided. In particular, boundary changes on auxiliary higher models do not alter the original compact-fiber section dimension when only exceptional terms have been introduced on that fiber.

\begin{lemma}[Finite pullback]\label{num:finite}
If $\pi:V'\to V$ is a dominant generically finite morphism of normal
projective varieties and $D$ is a rational Cartier divisor on $V$, then
\[
 \kappa(V',\pi^*D)=\kappa(V,D).
\]
\end{lemma}

\begin{proof}
Stein factorization reduces to finite $\pi$, since a proper birational
morphism onto a normal variety has direct image $\mathcal O_V$.
Clear denominators. Each homogeneous section upstairs is integral over
the downstairs section ring: its monic characteristic polynomial in
the finite function-field extension has coefficients in the appropriate
powers of $D$. Trivializing $D$ at each prime of $V$ shows that these
coefficients are regular there, hence everywhere by normality. The
two graded rings therefore have the same transcendence degree. If the
downstairs ring has no positive-degree section, the norm excludes one
upstairs. These observations prove the assertion in all cases.
\end{proof}

\begin{lemma}[Fields, sections, and specialization]
\label{prelim:generic}\label{lem:field-transfer}
For a proper variety and a line bundle, global sections and their
multiplication maps commute with extension of the ground field.
Consequently Iitaka dimension is unchanged by field extension.
For a finite surjective morphism \(\pi:Y\to X\) of normal projective
varieties and a rational Cartier divisor \(L\) on \(X\),
\[
 \kappa(Y,\pi^*L)=\kappa(X,L).
\]
Finite collections of morphisms and birational identifications spread
over fields of finite type. Section-system statements on geometric
generic fibers may be checked on very general complex fibers,
simultaneously in countably many divisible degrees, and conversely.

If \(b\subseteq b'\) are algebraically closed fields and \(I\) is a
geometrically integral \(b\)-variety, every nonempty open subset of
\(I_{b'}\) contains a point of \(I(b)\).
\end{lemma}

\begin{proof}
Proper cohomology and flat base change identify the section spaces,
compatibly with multiplication, and hence preserve the dimensions of
the section maps. The finite-morphism assertion follows from Lemma~\ref{num:finite}. Morphisms, rational maps, and the identities making two
rational maps inverse use finitely many coefficients and spread after
finitely many denominators are inverted. On the smooth proper locus, in each fixed degree we shrink to an open
where the section dimension is constant; cohomology and base change
then identify the complete fiber space. Spreading the section maps
and excluding the countably many proper closed complements gives the
very-general-fiber assertion.

For the last assertion, work on an affine open of \(I\). A regular
function after extension to \(b'\) is a finite sum
\(\sum_j c_jf_j\), where the \(c_j\in b'\) are linearly independent
over \(b\) and \(f_j\in\cO(I)\). If it vanishes on every \(b\)-point,
then every \(f_j\) does too and is zero. Thus \(I(b)\) remains Zariski
dense after field extension, proving the assertion.
\end{proof}

\begin{lemma}[Upper addition and interpolation]\label{prelim:addition}
For a projective fibration \(u:T\to R\) and a rational divisor \(D\) with nonempty divisible section systems, the dimension of its section-system image is at most the dimension of \(R\) plus the dimension of the restricted section-system image on the generic fiber. In particular it is at most
\[
\dim R+\kappa(T_{\overline\eta},D|_{T_{\overline\eta}}).
\]
If \(D\) is rationally effective and \(b\geq1\) is rational, then for every rational divisor \(M\),
\[
\kappa(D+M)\geq\kappa(D+bM).
\]
\end{lemma}

\begin{proof}
For a sufficiently divisible system, compare its image with the base using the graph of the two rational maps. The fiber image has dimension bounded by the restricted complete system; the fiber dimension theorem gives the assertion, and taking the maximum over divisible systems proves the first inequality. This is the usual upper-addition estimate; see \cite[Theorem~4, p.~357]{IitakaOriginal}.

For the second claim use
\[
D+M=b^{-1}(D+bM)+(1-b^{-1})D.
\]
After clearing denominators, multiplication by a fixed effective section of the second summand embeds the corresponding section systems of the first. Their ratios are unchanged. Taking their image dimensions proves the assertion.
\end{proof}

We use the Iitaka fibration theorem for a divisor with nonnegative section dimension \cite[Theorems~3--4]{IitakaOriginal}: on a suitable smooth model its connected generic fiber has restricted divisor of Kodaira dimension zero, and the base has dimension equal to the original section dimension. Its application to a log adjoint divisor is explained in the two-base construction, including the comparison under subsequent resolutions.

For analytic computations, we normalize
\(\ddc=\frac{i}{2\pi}\partial\bar\partial\), so that
\(\ddc\log|s_D|^2=[D]\) in a local holomorphic frame. A weight is a local potential for a singular Hermitian metric; differences of weights on the same line bundle are globally defined functions. Smooth reference metrics are understood whenever a potential or a norm is compared globally. Positive universal normalization constants do not affect the integral estimates below.



Let $f:U\to V$ be the morphism in Theorem~\ref{thm:main}, and put
$d=\kappa(F)$.  If $d=-\infty$, the asserted inequality follows from the
convention for $-\infty$.  We therefore assume throughout this section that
$d\geq 0$.  We first obtain the contribution of $\kbar(V)$ and then construct
the bases used to control variation.  The subadditivity used in this section is
Theorem~\ref{spec:logarithmic}.

\subsection{Compactification and the first Iitaka base}

\begin{lemma}\label{it:compactification}
There are smooth connected projective compactifications $X$ and $Y$ of $U$
and $V$, respectively, and a surjective morphism $X\to Y$ with connected
fibers extending $f$, such that the reduced divisors
\[
 D_X=X\setminus U,\qquad D_Y=Y\setminus V
\]
have simple normal crossings and the inverse image of $V$ is exactly $U$.
In particular, $D_X$ has no component dominating $Y$, and
\begin{equation}\label{it:base-bound}
 \kappa(X,K_X+D_X)=\kbar(U)
 \geq d+\kbar(V).
\end{equation}
\end{lemma}

\begin{proof}
Take a smooth projective compactification of $V$ and resolve its boundary.
Since $f$ is projective, embed $U$ as a closed subscheme of a projective
space over $V$ and take its closure over the resulting projective base $Y$.
Equivalently, one can take the closure of the graph of $f$ in projective
compactifications.  Over $V$ this closure is already $U$: the original
family is proper over $V$, so no additional limit point is needed above
that open set.  Resolve the closure and its boundary, with all centers
outside the smooth open $U$.  Resolution and elimination of indeterminacy
in characteristic zero give the required morphism between smooth projective
varieties and the SNC boundary; see \cite{BierstoneMilmanFunctorial}.

The resulting morphism has a geometrically integral generic fiber.  Indeed,
the original proper morphism has connected fibers and normal base, so its
Stein factorization gives $f_*\cO_U=\cO_V$.  Thus $\C(V)$ is relatively
algebraically closed in $\C(U)$; in characteristic zero this extension is
regular.  The finite part of the Stein factorization of the compactified
morphism is consequently birational over $Y$, and hence is an isomorphism
because $Y$ is normal.  This proves connectedness of its fibers.

As $f^{-1}(V)=U$, every point of $D_X$ lies over $D_Y$, and
$\Supp(f^*D_Y)\subseteq\Supp(D_X)$.  A very general fiber lies over $V$,
has no boundary, and has ordinary Kodaira dimension $d$.
Theorem~\ref{spec:logarithmic} now applies to these smooth projective varieties,
connected fibers, and reduced SNC boundaries.  Its conclusion is
\eqref{it:base-bound}, by the compactification definition of logarithmic
Kodaira dimension.
\end{proof}

Fix, once and for all, a nonzero section
\begin{equation}\label{it:base-form}
 \xi\in H^0\bigl(Y,m(K_Y+D_Y)\bigr),\qquad m>0.
\end{equation}
Such a section exists because $\kbar(V)\geq0$.  It is essential here that
$\xi$ is logarithmic: no nonnegative ordinary Kodaira dimension of $Y$ is
assumed.

We recall the boundary convention for all subsequent source models.
If $\pi:X_1\to X$ is a smooth birational model, set
\[
 D_1=\pi_*^{-1}D_X+\operatorname{Exc}(\pi)_{\mathrm{red}}.
\]
The SNC discrepancy formula gives
\begin{equation}\label{it:exceptional-sections}
 K_{X_1}+D_1
 =\pi^*(K_X+D_X)+\sum_E A(E;X,D_X)E,
 \qquad A(E;X,D_X)\geq0,
\end{equation}
where the sum is exceptional.  For every positive integer $a$, divisorial
testing on the normal target therefore gives
\[
 \pi_*\cO_{X_1}\bigl(a(K_{X_1}+D_1)\bigr)
 =\cO_X\bigl(a(K_X+D_X)\bigr).
\]
Thus these modifications preserve the entire adjoint section ring, not
just its Iitaka dimension.  The same argument permits a pulled-back
rational twist in divisible degrees.  After each modification we resolve
the boundary again, using the same convention.

\begin{proposition}\label{it:diagram}
There are smooth connected projective varieties $X_*,X',W,Z$, a birational
morphism $X'\to X_*$, and morphisms
\[
 X'\xrightarrow{x}W\xrightarrow{g}Y,
 \qquad h:W\to Z,\qquad q=h\circ x,
\]
with the following properties.
\begin{enumerate}[label=\textup{(\roman*)}]
 \item The variety $X_*$ is a smooth model of $X$ carrying a resolved
 absolute Iitaka map for $K_X+D_X$.  Its boundary $D_*$, and the boundary
 $D'$ on $X'$, follow the strict-transform-plus-reduced-exceptional
 convention.  The map $q$ is the pullback of this Iitaka map, and
 \begin{equation}\label{it:iitaka-dimension}
  \dim Z=\kbar(U).
 \end{equation}
 \item The morphisms $x,g,h$ are surjective and have connected fibers and
 geometrically integral generic fibers.  The induced map from $W$ to
 its image in $Y\times Z$ is generically finite.  In particular,
 $W_z\to g(W_z)$ is generically finite for very general $z$.
 \item The divisor
 $D_W^0=(g^*D_Y)_{\mathrm{red}}$ is SNC.  On very general fibers over $Z$,
 all the relevant varieties and boundary strata are smooth, and
 \begin{equation}\label{it:absolute-fiber}
  \kappa\bigl(X_{*,z},K_{X_{*,z}}+D_*|_{X_{*,z}}\bigr)
  =\kappa\bigl(X'_z,K_{X'_z}+D'|_{X'_z}\bigr)=0.
 \end{equation}
\end{enumerate}
One may make further compatible smooth birational modifications, retaining
these assertions.  After the finite preparations in this section, $X_*$
can be fixed as an earlier model, and every subsequent source model is
taken to dominate it.
\end{proposition}

\begin{proof}
By \eqref{it:base-bound}, $K_X+D_X$ has nonnegative section dimension.
The divisor Iitaka theorem gives a resolved map with base dimension
$\kbar(U)$ and with restriction of the divisor having section dimension
zero on a very general fiber.  The needed statement is
\cite[Theorem~3, p.~357, and its proof in \S4, pp.~361--363]{IitakaOriginal}.
The stabilized section field is relatively algebraically closed in
$\C(X)$ by \cite[\S3, Proposition~1, p.~359]{IitakaOriginal}.  In the case
$\kbar(U)=0$, the map is the map to a point.

For clarity, the fiber conclusion concerns the full restricted section
system.  On a resolution, a sufficiently divisible system has moving part
pulled back from a big system on the Iitaka base.  In a divisible degree
there is consequently a section after subtracting a pulled-back ample
divisor.  Any extra independent plurisections on the generic fiber extend
after an ample pullback twist.  Multiplying by the preceding section
would then enlarge the stabilized section field, which is impossible.
Doing this in each degree gives the assertion on very general fibers.
Formula~\eqref{it:exceptional-sections} allows the theorem to be applied
to the log divisor on each higher source model itself.  Adjunction on a
smooth fiber identifies its restriction with that fiber's log canonical
divisor.  This proves \eqref{it:absolute-fiber}, including after further
compatible resolutions and birational changes of the Iitaka base.

Let $K=\C(X)$, $L_Y=\C(Y)$, and $L_Z=\C(Z)$, viewed as subfields of $K$.
The combined image in $Y\times Z$ has function field $L_YL_Z$.  Let $L_W$
be the relative algebraic closure of this compositum in $K$, a finite
extension of $L_YL_Z$.  Normalize the combined image in $L_W$, resolve it,
and resolve the resulting rational map from a higher source model.
These operations give the displayed morphisms.

The field $L_W$ is relatively algebraically closed in $K$ by construction.
Moreover, $L_Y$ and $L_Z$ are relatively algebraically closed in $K$, so
each is relatively algebraically closed in $L_W$.  In characteristic zero
the three corresponding generic field extensions are regular.  This
proves geometric integrality of the generic fibers of $x,g,h$; properness
and the normality of their targets give connectedness by Stein
factorization.  Finiteness of $L_W/(L_YL_Z)$ gives generic finiteness to
the combined image and hence generic finiteness of the indicated fiber
maps.

Finally resolve $D_W^0$ and the source boundary and lift all the maps.
Outside a proper closed subset of $Z$, the morphisms and every dominating
boundary stratum are smooth, and the nondominating strata are absent.
This follows from generic smoothness applied to the finitely many strata.
It gives the asserted very general SNC fiber data.  The section-ring and
field arguments just given are unchanged by these birational operations.
\end{proof}

The two maps just constructed are shown in
Figure~\ref{fig:combined-base}. The lower arrows in the figure are
the next objective; they will identify the required subfield of
$\C(Y)$.
\begin{figure}[ht]
\centering
\begin{tikzcd}[row sep=large,column sep=huge]
 & X' \arrow[d,"x"] & \\
 & W \arrow[dl,"g"'] \arrow[dr,"h"] & \\
Y \arrow[dr,dashed] & & Z \arrow[dl] \\
 & T_0 &
\end{tikzcd}
\caption{The two maps from the total space: $g x$ is the original
fibration and $h x$ is the absolute log Iitaka map. The construction
of $W$ makes its map to $Y\times Z$ generically finite onto its image. The proof
then constructs $T_0$ by parametrizing $g(W_z)$ and shows that both
base fields contain $\C(T_0)$. The lower arrows are conclusions of
the argument, not assumptions. Resolving the dashed map later
gives $S\to T_0$ with $S$ birational to $Y$.}
\label{fig:combined-base}
\end{figure}


\subsection{Logarithmic restrictions to the Iitaka fibers}

Write
\[
 n=\dim Y,\qquad r=\dim W_z,\qquad p_1=n-r.
\]
Generic finiteness of $W_z\to g(W_z)$ gives $0\leq r\leq n$, so
$p_1\geq0$.  The next argument explains how the logarithmic form on $Y$
produces a form of the correct degree on $W_z$.

\begin{lemma}\label{it:contraction}
For very general $z\in Z$, the form $\xi$ in \eqref{it:base-form} induces
nonzero sections of $m(K_{W_z}+D_W^0|_{W_z})$ by contraction with suitable
$p_1$ tangent directions at $z$.  In particular,
\[
 \kappa\bigl(W_z,K_{W_z}+D_W^0|_{W_z}\bigr)\geq0.
\]
\end{lemma}

\begin{proof}
Work over an open subset of $Z$ on which $h$ and all the boundary strata
are smooth.  On its inverse image the logarithmic cotangent sequence is
exact:
\[
 0\longrightarrow h^*\Omega_Z^1
 \longrightarrow\Omega_W^1(\log D_W^0)
 \longrightarrow\Omega_{W/Z}^1(\log D_W^0)
 \longrightarrow0.
\]
The last bundle has rank $r$.  Since $n=r+p_1$, the exterior-power
filtration has no term of base degree below $p_1$, and its first quotient
is the canonical map
\begin{equation}\label{it:wedge-quotient}
 \bigwedge^n\Omega_W^1(\log D_W^0)
 \longrightarrow
 h^*\!\bigwedge^{p_1}\Omega_Z^1
 \otimes\det\Omega_{W/Z}^1(\log D_W^0).
\end{equation}
Pullback of logarithmic differentials under $g$ is regular here: locally,
the pullback of an equation for $D_Y$ is a monomial in equations for
$D_W^0$ times a unit.  Thus $g^*\xi$ is a section of the $m$th tensor
power of the left side of \eqref{it:wedge-quotient}.  Apply the quotient
tensorwise and, on $W_z$, evaluate on any $p_1$ fixed vectors in $T_zZ$.
The result is a global section of
\[
 \bigl(\det\Omega_{W/Z}^1(\log D_W^0)|_{W_z}\bigr)^{\otimes m}
 \simeq\cO_{W_z}\bigl(m(K_{W_z}+D_W^0|_{W_z})\bigr).
\]
This construction is independent of local lifts of the chosen tangent
vectors: replacing one lift by a relative tangent vector inserts $r+1$
relative vectors into an alternating form, so contributes zero.  It is
therefore valid along the entire boundary of the fiber, rather than
only on its interior.

At a generic point of $W_z$, the differential of $g$ has rank $n$, while
its restriction to the tangent space of $W_z$ has rank $r$.  The induced
normal map from $T_zZ$ consequently has rank $p_1$.  Some choice of the
$p_1$ vectors gives a nonzero determinant, and hence a nonzero contracted
section.  When $p_1=0$, the same construction simply uses the relative
top-degree quotient, with no tangent vectors to choose.
\end{proof}

\begin{lemma}\label{it:zerofiber}
Let $J$ be the smooth geometric generic fiber of $x$, or a simultaneous
very general complex fiber.  Then
\begin{equation}\label{it:zero-equations}
 \kappa\bigl(W_z,K_{W_z}+D_W^0|_{W_z}\bigr)=0,
 \qquad
 \kappa(J)=\kappa\bigl(J,K_J+D'|_J\bigr)=0.
\end{equation}
The source-fiber equalities \eqref{it:absolute-fiber} hold on the same
very general locus.
\end{lemma}

\begin{proof}
First, $\kappa(J)\geq0$.  To see this using the ordinary canonical
divisor, take a very general compact $Y$-fiber of $X'\to Y$.  It is a
smooth birational model of the original fiber $F$ and therefore carries
a nonzero ordinary pluricanonical section because $d\geq0$.  The
restriction of this section to a general fiber of its map to $W_y$ is
nonzero: its zero divisor cannot contain all those fibers.  Adjunction
identifies this restriction, up to the one-dimensional determinant of
the base cotangent space, with an ordinary pluricanonical section on
$J$.  Since $D'|_J$ is effective, its log canonical divisor also has
nonnegative section dimension.

Now restrict $x$ to a very general fiber over $Z$:
\[
 x_z:(X'_z,D'|_{X'_z})\longrightarrow
       (W_z,D_W^0|_{W_z}).
\]
The source and base are smooth connected projective varieties, the
fibers are connected, and both displayed boundaries are reduced SNC.
Every component of the inverse base boundary is in the source boundary,
because both lie over $D_Y$. Thus Theorem~\ref{spec:logarithmic}
applies. Its conclusion, together with
\eqref{it:absolute-fiber} and Lemma~\ref{it:contraction}, is
\[
 0\geq
 \kappa\bigl(J,K_J+D'|_J\bigr)
 +\kappa\bigl(W_z,K_{W_z}+D_W^0|_{W_z}\bigr)\geq0.
\]
Both terms vanish.  The already proved ordinary nonnegativity, and
monotonicity under addition of an effective divisor, give
$\kappa(J)=0$ as well.  The simultaneous geometric-generic and very
general interpretations follow by base change in the countably many
section degrees.
\end{proof}

Although the original compactification has no horizontal boundary, a
later exceptional divisor can be horizontal for $x$.  Thus the proof
does not assert $D'|_J=0$; it proves both dimensions in
\eqref{it:zero-equations}.  On a very general $Y$-fiber, however, every
exceptional divisor that meets that fiber remains exceptional over the
original $F$: its center has codimension at least two after restriction,
by the dimension formula for a center dominating $Y$.  Accordingly,
\eqref{it:exceptional-sections} on this compact fiber identifies its
logarithmic section dimension with its ordinary dimension $d$.

\subsection{The image family and the second base}

\begin{proposition}\label{it:dimensions}
The models of Proposition~\ref{it:diagram} may be chosen together with
smooth connected projective varieties $S,T_0$ and morphisms in the
commutative diagram
\begin{equation}\label{it:full-diagram}
\begin{tikzcd}[column sep=large,row sep=large]
 X' \arrow[r,"x"] \arrow[rr,bend left=28,"q=h\circ x"]
 & W \arrow[r,"h"] \arrow[d,"g_S"] \arrow[dl,"g"']
 & Z \arrow[d]
 \\
 Y
 & S \arrow[l,"\rho"] \arrow[r]
 & T_0
\end{tikzcd}
\end{equation}
such that $\rho$ is birational, $g=\rho\circ g_S$, and
\begin{equation}\label{it:product-map}
 W\longrightarrow(S\times_{T_0}Z)_{\mathrm{main}}
 \quad\text{is dominant and generically finite}.
\end{equation}
Both $Z\to T_0$ and $S\to T_0$ have geometrically integral generic
fibers.  With $D_S=(\rho^*D_Y)_{\mathrm{red}}$ made SNC, one has
\begin{equation}\label{it:dimension-equations}
 \begin{gathered}
  \dim(W/Y)=d,\qquad
  \dim T_0=p_1,\qquad
  \dim Z=d+\dim T_0,\\
  \kappa\bigl(S_t,K_{S_t}+D_S|_{S_t}\bigr)\geq0
 \end{gathered}
\end{equation}
for very general $t\in T_0$.
\end{proposition}

\begin{proof}
\emph{The relative dimension over $Y$.}
The dimension of $W_y$ is the image dimension of the absolute Iitaka
coordinates restricted to a very general $Y$-fiber; the finite extension
in the definition of $W$ does not change this dimension.  Those
coordinates form a subsystem of a log pluricanonical system on that
fiber.  By the exceptional section comparison just explained, their
image dimension is at most $d$.  Conversely, upper addition for the
ordinary canonical divisor on that compact fiber, whose general fiber
over $W_y$ is $J$, gives
\[
 d\leq\dim W_y+\kappa(J)=\dim W_y.
\]
Here we use the upper addition inequality for section dimensions with
irreducible general fibers, as in \cite[Theorem~4, p.~357]{IitakaOriginal}; no
nonnegativity or subadditivity hypothesis on $W_y$ is involved.
Consequently $\dim(W/Y)=d$.

\emph{The differential of the image family.}
Let $C_z=g(W_z)\subseteq Y$ with its reduced structure.  The scheme
image of $W\to Y\times Z$ has geometrically integral generic fiber
over $Z$: its fiber function field lies between $\C(Z)$ and the
regular extension $\C(W)$.  Generic flatness, followed by restriction
to a smaller open, therefore yields a flat family of reduced integral
subschemes $C_z$ of $Y$. Hilbert representability
\cite[Theorems~3.1--3.2 and the discussion following Proposition~3.8]{GrothendieckHilbert}
therefore gives a parameter map
\[
 \tau:Z\dashrightarrow\operatorname{Hilb}(Y).
\]
We claim that the dimension of its image is exactly $p_1$.

Fix a very general $z$ and a basis of $T_zZ$.  Perform the contraction
in Lemma~\ref{it:contraction} for every $p_1$-tuple of basis vectors.
By Lemma~\ref{it:zerofiber}, all nonzero resulting sections in degree
$m$ are proportional.  Indeed, two independent sections in one degree
would give a nonconstant ratio and positive Iitaka dimension.

At a general smooth point $y\in C_z$, the normal deformation map is
\begin{equation}\label{it:normal-map}
 T_zZ\longrightarrow N_{C_z/Y,y}.
\end{equation}
It is computed by lifting a base direction to $W$, applying $dg$, and
taking its class modulo $T_yC_z$.  This is independent of the lift and
has rank $p_1$, since $g$ is dominant and $W_z\to C_z$ is generically
finite.  Locally trivialize the tangential determinant and the form
$\xi$.  Each contracted section is the $m$th power of the corresponding
$p_1$-minor of \eqref{it:normal-map}, multiplied by a common nonzero
factor.  Proportionality of the sections says that the $m$th powers of
all ratios of nonzero minors are constant as the point of $W_z$ varies.
The ratios themselves are constant: a rational function whose positive
power is constant is algebraic over $\C$, and hence is constant on an
integral complex variety.  The Pl\"ucker coordinates of
\eqref{it:normal-map} thus have constant ratios.  Its kernel, a subspace
of the fixed vector space $T_zZ$, is consequently independent of the
general point $y$.

To pass from this pointwise calculation to the parameter map, write
$\mathcal I_z$ for the ideal of $C_z$ in $Y$.  A Hilbert tangent vector
is a homomorphism
\[
 \mathcal I_z/\mathcal I_z^2\longrightarrow\cO_{C_z}.
\]
Such a homomorphism is zero if it is zero at the generic point: the
target is the structure sheaf of an integral scheme and has no
torsion.  On the general smooth locus of $C_z$, its normal value is
precisely \eqref{it:normal-map}.  It follows that $d\tau_z$ has the
same kernel as that normal map and therefore has rank $p_1$.  In
characteristic zero, generic differential rank equals the dimension
of the reduced image.  This proves the claim.

This argument includes the extremes.  If $p_1=0$, every $C_z$ has
dimension $\dim Y$ and is therefore $Y$, so the family is constant.
If $r=0$, the smooth connected $W_z$ is a point and its parameter
map has rank $\dim Y=p_1$.

\emph{Construction of $S$ and $T_0$.}
Let the function field of $T_0$ be the relative algebraic closure of
the image parameter field in $\C(Z)$, and take a smooth projective
model.  Then $\dim T_0=p_1$, and $Z\to T_0$ has geometrically
integral generic fiber.  Pull the universal image family back to
$T_0$ and take the component dominated by $W$, with a smooth
projective resolution denoted by $S$.  After compatible resolutions
all the maps in \eqref{it:full-diagram} are morphisms.  Its right
square commutes, and the induced map \eqref{it:product-map} is
dominant and generically finite: over a general $z$ it is the
generically finite map to the corresponding image member.

The generic fiber of $S\to T_0$ has dimension $r$.  Thus
$\dim S=r+p_1=\dim Y$, and its dominant map $\rho:S\to Y$ is
generically finite.  In fact it is birational.  The maps give the
field inclusions
\begin{equation}\label{it:field-tower}
 \C(Y)\subseteq\C(S)\subseteq\C(W)\subseteq\C(X).
\end{equation}
The first extension is algebraic, whereas $\C(Y)$ is relatively
algebraically closed in $\C(X)$.  Therefore $\C(S)=\C(Y)$.
In particular, the construction embeds $\C(T_0)$ in the original
base field $\C(Y)$.

For completeness, geometric integrality of the generic fiber of
$S\to T_0$ can be checked in the same field tower.  The field
$\C(T_0)$ is relatively algebraically closed in $\C(Z)$ by
construction, and $\C(Z)$ is relatively algebraically closed in
$\C(W)$.  It follows that $\C(T_0)$ is relatively algebraically
closed in $\C(W)$, and hence in its subfield $\C(S)$.  The
extension $\C(S)/\C(T_0)$ is therefore regular.  This also verifies
the geometric connectedness independently of the universal-family
description.

\emph{Dimensions and the logarithmic base fiber.}
The equality already proved over $Y$ gives
\[
 \dim W=\dim Y+d=\dim Z+r.
\]
Together with $\dim T_0=\dim Y-r$, this yields
$\dim Z=d+\dim T_0$.

Resolve the support of $D_S=(\rho^*D_Y)_{\mathrm{red}}$, lifting the
other maps once more.  Logarithmic pullback gives a nonzero section
$\rho^*\xi$ of $m(K_S+D_S)$.  Restrict it to a very general smooth
fiber $S_t$; this restriction is nonzero, since its zero divisor
cannot contain all fibers.  Adjunction, with the determinant of
$T_t^*T_0$ a constant one-dimensional factor, gives a nonzero
section of $m(K_{S_t}+D_S|_{S_t})$.  This proves the last assertion
of \eqref{it:dimension-equations}.  All these final resolutions
preserve the section and fiber properties in
Proposition~\ref{it:diagram}.  We now fix an earlier resolved
Iitaka source $X_*$ and take the common source $X'$ above it;
later source modifications will continue to dominate this fixed
model.
\end{proof}

\begin{remark}\label{it:constant-field-interface}
The field identifications in \eqref{it:field-tower} will also control
the constant field in the final descent.  Inside the prescribed
algebraic closure $\Omega$ of $\C(Y)$, set
$b=\overline{\C(T_0)}$.  The regular extensions
$\C(Y)/\C(T_0)$ and $\C(W)/\C(Y)$ remain regular after the
corresponding base changes.  Thus, if $I$ and $H_0$ are the
geometric generic fibers of $S$ and $W$ over $T_0$, their map
$H_0\to I$ has geometrically integral generic fiber.  Moreover,
$b(I)=b\C(Y)\subseteq\Omega$ is algebraic over $\C(Y)$, so
$\Omega$ is also an algebraic closure of $b(I)$.  These are
statements about the fields in the original diagram; they do not
assert constancy of the $x$-fiber or of the whole original fiber.
\end{remark}
